
\section{Compute, Artifacts, and Reproducibility}
\label{sec:appendix:artifacts}

\subsection{Compute}
\label{sec:appendix:compute}

The campaign's canonical run cost \CostCalls{} model calls, \CostTokens{} tokens and
\CostModelHours{} model-hours across both roles, all \LabSeeds{} seeds and all three arms; wall
clock comes in lower, because the run fans out across workers. Call counts are identical across
arms, since every arm meets the same worlds, the same horizon and the same frozen task stream --- an
arm difference there would mean the arms were never matched and no contrast in \autoref{sec:lab}
would be licensed, so we assert it rather than describe it. The manipulation lives instead in
maintainer input tokens, which rise with what each arm's handoff carries
(\CostMaintainerInRange{} from the uncommented arm to the commented one). Both roles run
\texttt{claude-haiku-4-5}; each world's one-line objective is generated once at build time by
\texttt{claude-sonnet-5} and then frozen with the world, identically for every arm.

We train no model, so our budget is inference and parsing rather than gradient steps. Both halves
run on CPU-only cloud containers (8 vCPU, 16\,GiB, no accelerator). The corpus half is network- and
parse-bound, cloning \FieldRepos{} repositories bloblessly and segmenting every historical version
of every context file, and we enforce its wall-clock budget in code rather than watching it by
hand. The controlled experiment issues those calls against a hosted API, so the
parameter counts of \texttt{claude-haiku-4-5} and \texttt{claude-sonnet-5} are not public and we
cannot report them. We do pin the model versions and every decoding setting that varies
(\texttt{max\_tokens} per role, extended thinking disabled for objective generation) in the run's
resolved config. We otherwise decode at API defaults and ran no hyperparameter search over them.
The design's knobs are the feedback conditions of \apxref{sec:appendix:regime} and the comment
protocol itself, and we report both with their failures rather than only their survivor
(\autoref{tab:ablations}).

\subsection{What we consume, and under what license}
\label{sec:appendix:licenses}

We consume two artifacts, and we use each for
research, the purpose each was released for. The agent-context sampling frame
is the repository selection of \citet{chatlatanagulchai2025agentreadmes},
distributed as a replication package that carries no license file. We therefore
take from it only the list of repository URLs and re-derive every byte of
content from the public histories of those repositories themselves, which is
also the measurement this study requires
(\apxref{sec:appendix:construction}). We build the controlled experiment's
worlds from the IFEval input set \citep{zhou2023ifeval}, released under
Apache~2.0 as part of \texttt{google-research}. We redistribute no repository
content. We release only derived per-instruction tables and the code that
produces them from the public sources above.

\subsection{What we reimplement}
\label{sec:appendix:reimplementation}

We implemented the \LabVerifierTypes{}
constraint verifiers of \autoref{tab:verifiers} ourselves rather than reuse
IFEval's evaluation code, because our testbed inverts the benchmark: a verifier
must run against a maintained prompt and return a censored complaint rather
than score a completion (\apxref{sec:appendix:regime}). A reimplementation can
disagree with the original at the margin. Every satisfaction rate in this paper
is therefore a rate against our verifiers, comparable across arms rather than
against published IFEval numbers.

\subsection{Packages that determine a reported number}
\label{sec:appendix:packages}

We match across versions
with \texttt{rapidfuzz} at a similarity threshold of \MatchFuzzyThreshold{}
(\apxref{sec:appendix:tracking}), and we estimate the deletion hazard with the
Nelson--Aalen estimator from \texttt{lifelines}. We fit the frailty and
interaction models with our own piecewise-exponential EM
(\apxref{sec:appendix:frailty}) rather than a library routine, so this appendix
specifies that estimator instead of referring to one. We call models through
the Anthropic Python SDK. A lockfile resolves exact versions and freezes them
into each run's \texttt{env.lock}, so a reader reproducing a number resolves the
same dependency set by construction rather than by matching a printed version
string.

\subsection{Coverage, and what the corpus is not documented for}
\label{sec:appendix:coverage}

The corpus is
\FieldRepos{} public GitHub repositories carrying \texttt{CLAUDE.md},
\texttt{AGENTS.md}, or \texttt{copilot-instructions.md}, yielding
\FieldMultiVersionFiles{} multi-version files, \FieldMatches{} tracked
transitions, and \FieldSpells{} instruction spells, of which
\FieldDeletions{} end in a deletion (\apxref{sec:appendix:construction}). We fit
nothing to held-out data, so we run no train/test split: we estimate on the
whole corpus, and the controlled experiment draws its worlds by seeded
permutation of the eligible IFEval pool (\apxref{sec:appendix:worlds}). We did
not measure the corpus's natural language distribution, its programming
language distribution, its domain composition, or any demographic property of
its authors. No result here should be read as covering non-English
instructions. Limitations states this as a bound on our scope rather than a gap
to be filled by assumption.

\subsection{Authorship metadata}
\label{sec:appendix:authorship}

Commit metadata identifies people, and we
reduce it at extraction to one bit per file: whether at least two distinct
non-bot authors have touched it, the covariate of
\apxref{sec:appendix:interaction}. Nothing downstream of that reduction
carries a name or an address, and no released artifact does.
